\subsection{Enfoque y organización del Diseño}

El sistema propuesto consiste en una base de datos orientada a la gestión y análisis de información vinculada al ámbito del básquetbol, incluyendo entidades como clubes, jugadores, competencias e inscripciones. El diseño se basa en un modelo entidad–relación (DER) previamente definido, el cual ha sido transformado a un modelo relacional que garantiza integridad, consistencia y eficiencia en el acceso a los datos.

La organización del diseño sigue un enfoque modular, donde cada entidad representa un concepto específico del dominio y se vincula mediante relaciones claramente definidas. Se priorizó la normalización de los datos para evitar redundancias y asegurar la coherencia de la información, facilitando su mantenimiento y escalabilidad a futuro.

\subsection{Dependencias e Interacciones}

La base de datos ha sido concebida como el núcleo central del sistema, pudiendo operar de manera autónoma o integrarse con aplicaciones externas, tales como sistemas de carga de estadísticas, plataformas de gestión deportiva o interfaces de usuario desarrolladas para entrenadores y organizadores.

\subsection{Antecedentes del Proyecto}

El desarrollo de esta base de datos surge de la necesidad de contar con una herramienta estructurada para la gestión de información deportiva y el análisis estadístico en el contexto del básquetbol formativo. En particular, se busca mejorar la organización de los datos recolectados durante competencias, así como facilitar su posterior análisis para la toma de decisiones.

El proyecto se enmarca dentro de un proceso más amplio de desarrollo de software, que incluye etapas de relevamiento de requerimientos, modelado conceptual (DER), diseño lógico y futura implementación. La experiencia previa en la recopilación manual o poco estructurada de datos evidenció la necesidad de un sistema más robusto, motivando la creación del presente diseño.